% from ega1-6-fr.tex
\oldpage[I]{146}
chacun des $f^{-1}(V_i)$ admet un recouvrement fini par des ouverts affines
$W_{ij}$ tels que $A(W_{ij})$ soit une algèbre de type fini sur $A(V_i)$, et
par suite aussi une algèbre de type fini sur $A(U)$~; d'où la conclusion.

Enfin, pour démontrer (iv), on peut se borner au cas où $S=Y$~; en effet,
$f_{(S')}$ est aussi égal à $f_{(Y_{(S')})}$, $f$ étant considéré comme un
$Y$-morphisme, et l'extension de la base étant $Y_{(S')}\to Y$
(\hyperref[I.3.3.9-fr]{3.3.9}). Soient alors $p$, $q$ les projections
$X_{(S')}\to X$ et $X_{(S')}\to S'$. Soit $V$ un ouvert affine dans $S$~;
$f^{-1}(V)$ est réunion finie d'ouverts affines $W_i$ dont chacun est tel que
$A(W_i)$ soit une algèbre de type fini sur $A(V)$
(\hyperref[I.6.3.2-fr]{6.3.2}). Soit $V'$ un ouvert affine de $S'$ contenu
dans $g^{-1}(V)$~; comme $f\circ p=g\circ q$, $q^{-1}(V')$ est contenu dans la
réunion des $p^{-1}(W_i)$~; d'autre part, l'intersection
$p^{-1}(W_i)\cap q^{-1}(V')$ s'identifie au produit $W_i\times_V V'$
(\hyperref[I.3.2.7-fr]{3.2.7}), qui est un schéma affine d'anneau isomorphe à
$A(W_i)\otimes_{A(V)}A(V')$ (\hyperref[I.3.2.2-fr]{3.2.2})~; ce dernier étant
par hypothèse une algèbre de type fini sur $A(V')$, la proposition est
démontrée.

\begin{corollary}[6.3.5]
\label{I.6.3.5-fr}
Soit $f:X\to Y$ un morphisme d'immersion. Si l'espace sous-jacent à $Y$
(resp. $X$) est localement noethérien (resp. noethérien), $f$ est de type fini.
\end{corollary}

On peut toujours supposer $Y$ affine (\hyperref[I.6.3.2-fr]{6.3.2})~; si
l'espace sous-jacent à $Y$ est localement noethérien, on peut en outre le
supposer noethérien et alors l'espace sous-jacent à $X$, qui en est un
sous-espace, est noethérien. Autrement dit, on peut supposer $Y$ affine et
l'espace sous-jacent à $X$ noethérien~; il existe alors un recouvrement de $X$
par un nombre fini d'ouverts affines $D(g_i)\subset Y$, où $g_i\in A(Y)$, tels
que $X\cap D(g_i)$ soit fermé dans $D(g_i)$ (donc un schéma affine
(\hyperref[I.4.2.3-fr]{4.2.3})), puisque $X$ est localement fermé dans $Y$
(\hyperref[I.4.1.3-fr]{4.1.3}). Alors $A(X\cap D(g_i))$ est une algèbre de type
fini sur $A(D(g_i))$, d'après (\hyperref[I.6.3.4-fr]{6.3.4, (i)}) et
(\hyperref[I.6.3.3-fr]{6.3.3}), et
$A(D(g_i))=A(Y)_{g_i}=A(Y)[1/g_i]$ est de type fini sur $A(Y)$, ce qui achève
la démonstration.

\begin{corollary}[6.3.6]
\label{I.6.3.6-fr}
Soient $f:X\to Y$, $g:Y\to Z$ deux morphismes. Si $g\circ f$ est de type fini,
et si $X$ est noethérien, ou $X\times_Z Y$ localement noethérien, $f$ est de
type fini.
\end{corollary}

Cela résulte aussitôt de la démonstration de
(\hyperref[I.5.5.12-fr]{5.5.12}) et de (\hyperref[I.6.3.5-fr]{6.3.5}) appliquée
au morphisme d'immersion $\Gamma_f$.

\begin{proposition}[6.3.7]
\label{I.6.3.7-fr}
Soit $f:X\to Y$ un morphisme de type fini~; si $Y$ est noethérien (resp.
localement noethérien), $X$ est noethérien (resp. localement noethérien).
\end{proposition}

On peut se borner à faire la démonstration lorsque $Y$ est noethérien. Alors $Y$
est réunion finie d'ouverts affines $V_i$ tels que les $A(V_i)$ soient des
anneaux noethériens. En vertu de (\hyperref[I.6.3.2-fr]{6.3.2}), chacun des
$f^{-1}(V_i)$ est réunion d'un nombre fini d'ouverts affines $W_{ij}$ tels que
les $A(W_{ij})$ soient des algèbres de type fini sur $A(V_i)$, donc des anneaux
noethériens~; cela prouve que $X$ est noethérien.

\begin{corollary}[6.3.8]
\label{I.6.3.8-fr}
Soit $X$ un préschéma de type fini sur $S$. Pour toute extension de la base
$S'\to S$ telle que $S'$ soit noethérien (resp. localement noethérien),
$X_{(S')}$ est noethérien (resp. localement noethérien).
\end{corollary}

Cela résulte de (\hyperref[I.6.3.7-fr]{6.3.7}), $X_{(S')}$ étant de type fini
sur $S'$ en vertu de (\hyperref[I.6.3.4-fr]{6.3.4, (iv)}).

On peut encore dire que dans un produit $X\times_S Y$ de $S$-préschémas, si
\emph{l'un} des
